% \setcounter{chapter}{\Session-2}
\renewcommand{\chaptername}{سوال}
\chapter{سوال ۲: Flow Matching برای سری‌های زمانی}

در این بخش نتایج پیاده‌سازی مدل \grayBox{Flow Matching} برای تولید سری‌های زمانی
\lr{log-return}
نماد \lr{SPY}
گزارش و تحلیل می‌شود. تمامی شکل‌ها از نوت‌بوک کد استخراج شده‌اند.

% ==========================
% بخش ۰: پرسش‌های نظری (Theory)
% ==========================
\section{پرسش‌های نظری}

\subsection*{سوال ۱: میدان برداری}
	extbf{الف) نقش میدان برداری:}
در Flow Matching، میدان برداری $v_\theta(x,t)$ «سرعت لحظه‌ای» حرکت جرم احتمال را تعیین می‌کند؛ یعنی اگر نمونه‌ی اولیه $x_0$ از توزیع ساده (مثلاً نویز گاوسی) در $t=0$ داشته باشیم، با حل معادله دیفرانسیل
$\frac{dx}{dt}=v_\theta(x,t)$
نمونه در طول زمان به سمت توزیع داده حرکت کرده و در $t=1$ به نمونه‌ای از توزیع هدف تبدیل می‌شود.

	extbf{ب) تفاوت/شباهت با Diffusion:}
Diffusion ذاتاً تصادفی است (زنجیره مارکوف یا SDE) اما Flow Matching یک مدل تعیین‌گر مبتنی بر ODE است. شباهت مهم این است که فرم‌های تعیین‌گر diffusion (مثل \lr{probability flow ODE} یا \lr{DDIM}) از نظر ایده به یادگیری یک مسیر ODE نزدیک می‌شوند.

\subsection*{سوال ۲: مسیرهای خطی}
	extbf{الف) چرا مسیر خطی؟ آیا همیشه بهینه است؟}
مسیر خطی $x_t=(1-t)x_0+tx_1$ ساده‌ترین مسیر بین دو نقطه است و سرعت هدف در این مسیر ثابت می‌شود: $v_{\text{target}}=x_1-x_0$؛ این سادگی یادگیری را آسان می‌کند. بهینگی «کلی» نیازمند کوپلینگ بهینه (\lr{Optimal Transport}) بین $x_0$ و $x_1$ است؛ در غیر این صورت مسیر خطی صرفاً یک انتخاب ساده و پایدار است.

	extbf{ب) اثر بر پایداری و هزینه:}
هدف رگرسیون ثابت و کراندار معمولاً پایداری آموزش را بهتر می‌کند و تابع هزینه به یک \lr{Least Squares} ساده تبدیل می‌شود.

\subsection*{سوال ۳: بازنویسی تابع هزینه (Score Matching)}
برای کرنل اغتشاش گاوسی
$q(x_t|x_0)=\mathcal{N}(x_t;x_0,\sigma_t^2 I)$
داریم:
\begin{equation}
\nabla_{x_t}\log q(x_t|x_0)=\frac{x_0-x_t}{\sigma_t^2}.
\end{equation}
بنابراین آموزش مدل امتیاز $s_\theta(x_t,t)$ با کمینه‌سازی خطای مربعی نسبت به این گرادیان (طبق نتیجه‌ی \lr{Vincent et al.}) به فرم
$\mathbb{E}[\|s_\theta(x_t,t)-(x_0-x_t)/\sigma_t^2\|_2^2]$
قابل بازنویسی است.

% ==========================
% بخش ۱: تابع هزینه و منحنی‌های خطا
% ==========================
\section{تابع هزینه و رفتار خطا}

\subsection{تابع هزینه}
تابع هزینه‌ی استفاده‌شده مطابق رابطه‌ی زیر است که سرعت پیش‌بینی‌شده‌ی میدان برداری را به سرعت هدف
\lr{$v_{\text{target}} = x_1 - x_0$}
روی مسیر خطی \lr{$x_t = (1-t)x_0 + tx_1$} نزدیک می‌کند:
\begin{equation}
\mathcal{L} = \mathbb{E}_{x_0, x_1, t} \left[ \left\| v_\theta(x_t, t) - (x_1 - x_0) \right\|_2^2 \right].
\end{equation}

\noindent\textbf{کد اصلی تابع هزینه در پیاده‌سازی:}
\begin{minted}[fontsize=\footnotesize]{python}
t = torch.rand(batch_size, device=device)  # t ~ Uniform(0,1)
t_view = t.view(batch_size, 1, 1)
x_t = (1.0 - t_view) * x0 + t_view * x1
v_target = x1 - x0
v_pred = model(x_t, t)
loss = F.mse_loss(v_pred, v_target)
\end{minted}

\subsection{منحنی خطای آموزش و آزمون}

\begin{figure}[H]
	\centering
	\maybeincludegraphics{width=0.8\linewidth}{../En_report/figures/fm_loss_curve.png}
	\caption{منحنی خطای آموزش و آزمون مدل Flow Matching در طول دوره‌های آموزش.}
	\label{fig:fm-loss}
\end{figure}

\noindent\textbf{توضیح شکل~\ref{fig:fm-loss}:}
کاهش یکنواخت خطای آموزش و نزدیکی خطای آزمون به آن نشان می‌دهد که مدل نه دچار بیش‌برازش شدید شده و نه کم‌برازش جدی دارد؛ فاصله‌ی زیاد بین دو منحنی می‌توانست نشانه‌ی بیش‌برازش باشد.

% ==========================
% بخش ۲: نمونه‌گیری و حل‌گرهای ODE
% ==========================
\section{نمونه‌های تولیدی و مقایسه‌ی حل‌گرها}

\subsection{گام ۳: نمونه‌گیری با حل‌گرهای ODE (Euler/Heun)}
در نمونه‌گیری، ODE به صورت عددی حل می‌شود. افزایش تعداد گام‌ها هزینه محاسبات را بالا می‌برد اما معمولاً کیفیت را بهتر می‌کند.
\begin{minted}[fontsize=\footnotesize]{python}
def euler_step(x, t, dt):
	return x + dt * v_theta(x, t)

def heun_step(x, t, dt):
	k1 = v_theta(x, t)
	x_e = x + dt * k1
	k2 = v_theta(x_e, t + dt)
	return x + 0.5 * dt * (k1 + k2)
\end{minted}

\subsection{نمونه‌های تولیدی (نرمال‌شده و بازنرمال‌شده)}

\begin{figure}[H]
	\centering
	\maybeincludegraphics{width=0.9\linewidth}{../En_report/figures/fm_generated_samples.png}
	\caption{نمونه‌های تولیدی مدل Flow Matching؛ ردیف بالا سری‌های زمانی نرمال‌شده و ردیف پایین همان نمونه‌ها بعد از بازنرمال‌سازی به مقادیر \lr{log-return} واقعی را نشان می‌دهد.}
	\label{fig:fm-gen}
\end{figure}

\noindent\textbf{توضیح شکل~\ref{fig:fm-gen}:}
دامنه‌ی نوسانات و خوشه‌بندی نوسانات (volatility clustering) در نسخه‌ی بازنرمال‌شده شبیه رفتار سری زمانی واقعی است و نشان می‌دهد مدل تا حد مناسبی ساختار آماری داده را یاد گرفته است.

\subsection{مقایسه‌ی کیفی داده‌ی واقعی و تولیدی}

\begin{figure}[H]
	\centering
	\maybeincludegraphics{width=0.9\linewidth}{../En_report/figures/fm_real_vs_gen.png}
	\caption{مقایسه‌ی کیفی چند نمونه‌ی واقعی (ردیف بالا) و نمونه‌های تولیدی متناظر (ردیف میانی و پایین).}
	\label{fig:fm-real-vs-gen}
\end{figure}

\noindent\textbf{توضیح شکل~\ref{fig:fm-real-vs-gen}:}
در بسیاری از نمونه‌ها، شکل کلی، دامنه و فرکانس نوسانات نمونه‌های تولیدی مشابه سری‌های واقعی است؛ انحراف‌های چشم‌گیر می‌تواند نشان‌دهنده‌ی ظرفیت ناکافی مدل یا نیاز به آموزش طولانی‌تر باشد.

% ==========================
% بخش ۳: مقایسه‌ی توزیع‌ها
% ==========================
\section{مقایسه‌ی توزیع‌های حاشیه‌ای}

\begin{figure}[H]
	\centering
	\maybeincludegraphics{width=0.9\linewidth}{../En_report/figures/fm_dist_comparison.png}
	\caption{مقایسه‌ی توزیع \lr{log-return} واقعی و تولیدی با استفاده از هیستوگرام، \lr{KDE}، نمودار \lr{Q-Q} و \lr{CDF}.}
	\label{fig:fm-dist}
\end{figure}

\noindent\textbf{توضیح شکل~\ref{fig:fm-dist}:}
هم‌پوشانی مناسب هیستوگرام و منحنی‌های \lr{KDE} نشان می‌دهد که مدل توزیع حاشیه‌ای بازده‌ها را تا حد خوبی تقریب زده است؛ اختلاف در دم‌ها (tail) بیان‌گر کم‌برآوردی یا بیش‌برآوردی رویدادهای نادر است که در نمودارهای \lr{Q-Q} و \lr{CDF} نیز قابل مشاهده است.

% ==========================
% بخش ۳.۵: ارزیابی آماری پایه
% ==========================
\section{ارزیابی آماری پایه}

\begin{table}[H]
    \centering
    \begin{tabular}{lcc}
        \toprule
        \textbf{شاخص} & \textbf{داده واقعی} & \textbf{داده تولیدی} \\
        \midrule
        میانگین & ... & ... \\
        واریانس & ... & ... \\
        نوسان‌پذیری & ... & ... \\
        \bottomrule
    \end{tabular}
    \caption{مقایسه شاخص‌های آماری پایه بین داده واقعی و تولیدی.}
\end{table}

\noindent\textbf{بحث در مورد نوسان‌پذیری:}
نوسان‌پذیری (مثلاً انحراف معیار بازده‌ها) یک شاخص کلیدی برای واقع‌گرایی مالی است. یک تولیدکننده خوب نوسان‌پذیری داده واقعی را در افق مجموعه‌داده مطابقت می‌دهد؛ نوسان‌پذیری به‌طور显著 پایین‌تر نشان‌دهنده فروپاشی حالت به سمت میانگین است، در حالی که نوسان‌پذیری به‌طور显著 بالاتر ممکن است نشان‌دهنده ناپایداری یا مقیاس‌بندی نادرست نمونه‌گیری باشد.

% ==========================
% بخش ۴: ارزیابی پیشرفته (SWD و خودهمبستگی)
% ==========================
\section{ارزیابی پیشرفته: فاصله Wasserstein بُرش‌خورده و خودهمبستگی}

\subsection{گام ۸: محاسبه SWD (طرح کلی)}
برای \lr{SWD}، داده‌ها روی چند جهت تصادفی پروجکت شده و میانگین فاصله Wasserstein تک‌بعدی گرفته می‌شود.
\begin{minted}[fontsize=\footnotesize]{python}
projections = torch.randn(n_proj, seq_len, device=device)
projections = projections / projections.norm(dim=1, keepdim=True)
real_proj = real @ projections.T
gen_proj = gen @ projections.T
swd = (torch.sort(real_proj, dim=0).values - torch.sort(gen_proj, dim=0).values).abs().mean()
\end{minted}

\begin{figure}[H]
	\centering
	\maybeincludegraphics{width=0.8\linewidth}{../En_report/figures/fm_autocorr.png}
	\caption{تحلیل خودهمبستگی در داده‌ی واقعی و تولیدی، به‌همراه جمع‌بندی شاخص‌هایی مانند \lr{SWD} و خطای میانگین مربعات خودهمبستگی.}
	\label{fig:fm-autocorr}
\end{figure}

\noindent\textbf{توضیح شکل~\ref{fig:fm-autocorr}:}
نمودار نوارهای خودهمبستگی بر حسب \lr{lag} نشان می‌دهد که مدل تا چه حد ساختار زمانی داده را حفظ کرده است؛ اگر میله‌های سبز (داده‌ی تولیدی) نزدیک به میله‌های آبی (داده‌ی واقعی) باشند، ساختار زمانی به‌خوبی یاد گرفته شده است. مقادیر کوچک \lr{SWD} و خطای خودهمبستگی نیز نشان‌دهنده‌ی نزدیکی توزیع و دینامیک مدل به داده‌ی واقعی است.

در مجموع، نتایج بالا نشان می‌دهد که مدل Flow Matching توانسته است هم توزیع حاشیه‌ای بازده‌ها و هم بخشی از وابستگی زمانی آن‌ها را به‌خوبی بازتولید کند، هرچند بهبودهای بیشتر (مثلاً افزایش ظرفیت مدل یا استفاده از حل‌گرهای دقیق‌تر \lr{ODE}) می‌تواند نتایج را بهتر کند.